\chapter{Le basi del Quantum Computing}

In tale capitolo sono presentate le ipotesi e le basi matematiche per affrontare quelli che sono i diversi argomenti che ci si troverà ad affrontare nei successivi capitoli

\section{Il primo postulato}
Il primo postulato della meccanica quantistica ci dice che: 

Ad ogni \textbf{sistema fisico} è associato uno \textbf{Spazio di Hilbert} \(\mathcal{H}\). Lo \textbf{stato} di tale sistema fisico è descritto da un vettore \(\psi\) di \(\mathcal{H}\) di norma unitaria.
Se due sistemi fisici, con spazi di Hilbert \(\mathcal{H}_1\) e \(\mathcal{H}_2\), vengono combinati in un unio sistema (facendoli interagire in qualche modo), lo spazio di Hilbert del sistema composto è dato dal prodotto tensoriale dei due spazi di Hilbert. \(\mathcal{H} \equiv \mathcal{H}_1 \otimes \mathcal{H}_2 \) \cite{giliberti2007limite}.

Per un sistema quantistico, i termini \textbf{isolato} e \textbf{chiuso} possono essere accostati, quindi saranno utilizzati come sinonimi, ed il significato puro è quello di sistema che non è accoppiato con un ambiente esterno non controllato.

Un sistema fisico quantistico è associato (come un sistema fisico) ad uno spazio di Hilbert, che sarà quindi chiamato spazio degli stati del sistema. Dato che parliamo (per ora) di sistemi isolati (non interagenti con un ambiente esterno non controllato), possiamo dire che il sistema è completamente descritto da un vettore di stato (per gli stati puri) con modulo unitario, a meno di una fase globale.
Quindi due vettori a modulo unitario, con solo una fase globale di differenza, sono visti come lo stesso stato puro.

\[
\ket{\psi} = e^{j \gamma} \ket{\phi}, \gamma \in [0,2\pi)
\]

Dove, considerando \(\mathbb{C}^n\) uno spazio vettoriale complesso (i cui elementi sono colonne), dove n rappresenta il numero di valori complessi totali per rappresentare il sistema nella sua interezza, allora:
\[
\ket{\psi}, \ket{\phi} \in \mathbb{C}^n
\]

Solitamente gli stati Quantistici possono andare a rappresentare diverse tipologie di caratteristiche che si vogliono analizzare nel sistema. Ad esempio lo spin degli elettroni all'interno degli atomi (possono essere o verso l'alto o verso il basso una volta osservati), oppure la polarizzazione orizzontale o verticale dei protoni. 
Il sistema più semplice in meccanica quantistia è il quantum bit (Qubit), dove lo spazio di stato è duo-dimensionale.

\begin{warn}
Il valore di n non deve essere confuso con il numero di Qubit. Dato che ogni Qubit sarà rappresentato da 2 valori complessi, la grandezza dello spazio sarà \(n=2^m\), dove \(m\) è il numero di Qubit del sistema.
\end{warn}

\subsection{Spazi di Hilbert e trasformazioni}
Uno spazio di Hilber è uno \textbf{Spazio Vettoriale Complesso} dotato di un prodotto interno (su cui si andrà a definire una norma), le cui successioni di vettori di cauchy convergono asintoticamente a un elemento dello spazio stesso. Tale convergenza, per spazi vettoriali a dimensione finita (come nei nostri casi), è sempre verificata. Tale spazio è spannato su una base ortonormale complessa finita. Le trasformazioni quantistiche saranno viste come matrici unitarie (\(U^\dagger U = U U^\dagger = I, U^\dagger = (U^T)^*\)) moltiplicate per i vettori di stato.

\begin{info}
Dato uno spazio vettoriale complesso \(\mathbb{V}\) dotato di prodotto scalare \(<-,->\), che ne definisce la norma: \(| \cdot | = \sqrt{<\cdot,\cdot>}\), e la distanza: \(d(V_1,V_2) = |V_1 - V_2|\).

Sì definisce successione di Cauchy, una successione di vettori: \(V_0, V_1, V_2, \dots\) tale per cui, \(\forall \epsilon \in \mathbb{R}, \epsilon > 0\)

\[
\exists N_0 \in \mathbb{N} : \forall m,n > N_0, d(V_m, V_n) \le \epsilon
\]

Uno spazio vettoriale prehilbertiano (dotato di prodotto interno) è detto \textbf{completo} se per ogni successione di cauchy (\(V_0, V_1, \dots\)), esiste un vettore \(V \in \mathbb{V}\) tale per cui:
\[
\lim_{n \to \infty} |V_n - V| = 0
\]
Quindi se tali successioni convergono in un elemento interno allo spazio stesso. Sì definisce, quindi, uno \textbf{Spazio di Hilbert}, uno spazio normato e completo.

Ogni spazio vettoriale complesso a dimensione finita, è automaticamente uno spazio completo. Quindi, ogni spazio vettoriale complesso normato a dimensione finita è uno spazio di Hilbert
\end{info}

\section{Qubit e Notazioni}

Un \textbf{Qubit} è il sistema quantistico più semplice, i cui stati possono essere descritti da uno spazio vettoriale complesso a due dimensioni. In generale, questo è dovuto al fatto che "la misurazione" starebbe nel capire in che verso sta spinnando l'elettrone in base ad una stessa "direzione". Ma questo è fuori dagli scopi del corso.

\subsection{Notazione di Dirac}
La \textbf{notazione di Dirac} (o \textbf{BraKet}) è un modo per simboleggiare i vettori che ci permette di "gestire" in maniera più coerente quello che è lo studio della computazione quantistica.
Considerando \( c_1, c_2, \dots, c_n \in \mathbb{C}\), allora si definiscono le seguenti notazioni.
Le notazioni di base sono quella della \textbf{Ket} (Vettore colonna):
\[
\ket{\varphi} =  
\begin{bmatrix}
c_1 \\
c_2 \\
\vdots \\
c_n    
\end{bmatrix}
\]

E della \textbf{Bra} (Vettore Riga), è il coniugato trasposto di \(\ket{\varphi}\):
\[
\bra{\varphi} = 
\begin{bmatrix}
c_1^* & c_2^* & \dots & c_n^*
\end{bmatrix}
\]

Date le precedenti notazioni, sì definisce il \textbf{prodotto interno} come:
\[
\bra{x}\ket{y} = \braket{x | y}
\]

Mentre si definisce il \textbf{prodotto esterno} come:
\[
\ket{x}\bra{y}
\]

\subsection{Qubit}
Nel caso del Qubit, essendo lo spazio di dimensione 2, si andrà a considerare quella che è la base canonica dello spazio, ovvero: \(\mathcal{B} = 
\left\{
\left[
\begin{smallmatrix}
1 \\
0    
\end{smallmatrix}
\right]
,
\left[
\begin{smallmatrix}
0 \\
1
\end{smallmatrix}
\right]
\right\}
\).
Tale base può essere anche scritta seguendo la seguente notazione:
\[
\ket{0} = \begin{bmatrix}1 \\ 0\end{bmatrix} 
\]
\[
\ket{1} = \begin{bmatrix}0 \\ 1\end{bmatrix}
\]

Pertanto possiamo scrivere la base del nostro spazio degli stati come:
\[
\mathcal{B} = \left\{ \ket{0}, \ket{1} \right\}
\]

Pertanto, uno stato quantistico può essere visto come la \textbf{superposizione} di stati, il cui risultato è un altro stato quantistico valido.
Quindi, uno stato \(\ket{\varphi} \in \mathbb{C}^2\) del Qubit, può essere visto come:
\[
\ket{\varphi} = \alpha_0 \ket{0} + \alpha_1 \ket{1}
\]

Dove, \(|\alpha_0|^2\) è la probabilità di ottenere esito \(\ket{0}\) dalla misura sulla base precedentemente presentata; mentre \(|\alpha_1|^2\) è la probabilità di misurare lo stato \(\ket{1}\). Dalla teoria della probabilità, quindi:
\[
\alpha_0, \alpha_1 \in \mathbb{C} : |\alpha_0|^2 + |\alpha_1|^2 = 1
\]
La somma dei moduli delle possibilità all'interno di un sistema quantistico dev'essere 1.

Altra base molto importante per la quantum computation è la \textbf{base di Hadamard}, ovvero:
\[
\mathcal{H} = \left\{ \ket{+},\ket{-}\right\}
\]

Dove sarà possibile esprimere la superposizione come:
\[
\ket{\varphi} = \beta_0 \ket{+} + \beta_1 \ket{-}
\]

Gli elementi della base di Hadamard possono essere espressi anche come elementi della base canonica, quindi è possibile andare a definire le seguenti relazioni:
\[
\ket{+} = \frac{1}{\sqrt{2}}\left(\ket{0} + \ket{1}\right)
\]
\[
\ket{-} = \frac{1}{\sqrt{2}}\left(\ket{0} - \ket{1}\right)
\]



\section{Multiple Qubits}
Quando si vanno a considerare sistemi con più qubit, allora entra in gioco il \textbf{prodotto tensoriale}.
Tale operazione si distingue dalla somma, prima di passare alla sua formalizzazione matematica è importante fare una considerazione sulla dimensione dello spazio vettoriale complesso in cui si andrà a lavorare. Con la somma tra due spazi vettoriali sì ottiene la seguente condizione:
\[
\mathbb{C}^n \oplus \mathbb{C}^m = \mathbb{C}^{n+m}
\]
Dove la base viene definira mediante l'unione delle due basi del sottospazio:
\[
\mathcal{B} = \{v_1,\dots,v_n, u_1, \dots, u_m\}
\]
Mentre con il prodotto tensoriale, lo spazio sarà di dimensione:
\[
\mathbb{C}^n \otimes \mathbb{C}^m = \mathbb{C}^{nm}
\]
Dove la base sarà costituita da:
\[
\mathcal{B} = \{v_1 \otimes u_1, v_1 \otimes u_2, \dots v_n \otimes u_m\}
\]

Se andiamo a considerare in maniera diretta due stazi: \(V\) e \(W\), associati a due qubit, questi saranno dotati di due basi ben distinte:
\[
\mathcal{B}_V = \{v_1, v_2\}
\]
\[
\mathcal{B}_W = \{w_1, w_2\}
\]

Dove, la base intesa per la somma degli spazi, sarà:
\[
\mathcal{B}_{V \oplus W}= \{v_1,v_2,w_1,w_2\}
\]

Mentre per il prodotto vettoriale, si avrà:
\[
\mathcal{B}_{V \otimes W} = \{v_1 \otimes w_1, v_1 \otimes w_2, v_2 \otimes w_1, v_2 \otimes w_2 \}
\]

Quando si va a considerare la base per il sistema di interazione tradue qubit, espressi nella base canonica, allora si avrà:
\[
\mathcal{B} = \{\ket{0} \otimes \ket{0}, \ket{0} \otimes \ket{1}, \ket{1} \otimes \ket{0}, \ket{1} \otimes \ket{1}\}
\]

Che più comunenemente, tale base, viene scritta come:
\[
\mathcal{B} = \{\ket{00}, \ket{01}, \ket{10}, \ket{11}\}
\]

Andando a fare un confronto con quello che è il caso classico, il numero di combinazioni con i qubit è uguale al numero di combinazioni possibili all'interno dei sistemi classici.

\begin{info}
Le posizioni all'interno del vettore con la ket, non perdono il loro significato funzionale. Quindi:
\[
\ket{00} = \ket{0_A 0_B}
\]
Dove \(A\) e \(B\), sono due qubit diversi e ben identificati. Il non specificare a quale sistema stiamo facendo riferimento è un abuso di notazione ( opiù una mancanza di notazione)
\end{info}

La differenza sostanziale sta nel fatto che nel caso classico, in un determinato istante, un registro può trovarsi in una sola delle sue combinazioni, mentre nel caso quantistico, una coppia di qubit ha una probabilità non nulla di potersi trovare in uno dei 4 possibili stati della base (dopo la lettura).


Detto questo, quindi, è possibile scrivere un qubit appartenente ad un sistema di 2 qubit come:
\[
\ket{\psi} 0 \alpha_0 \ket{00} + \alpha_1 \ket{01} + \alpha_2 \ket{10} + \alpha_3 \ket{11}
\]
\[
\alpha_0, \alpha_1, \alpha_2, \alpha_3 \in \mathbb{C}, |\alpha_0|^2 + |\alpha_1|^2 + |\alpha_2|^2 + |\alpha_3|^2 = 1
\]

Andando poi a generalizzare la cosa, per un sistema ad n-qubit, si avrà che:
\[
\ket{psi} = \sum_{i=0}^{2^n-1} \alpha_i \ket{i}
\]
\[
\alpha_i \in \mathbb{C} : \sum_{i}|\alpha_i|^2=1
\]

\section{Misurazione di un Qubit}
Finchè un qubit rimane in una sua superposizione, non vi è possibile accedere alle informazioni. Per farlo bisognerà "misurarlo". Il postulato che governa la misura di un qubit è il secondo postulato della meccanica quantistica. Tale postulato ci permette di dire che quando si va ad effettuare una misurazione su di un qubit è importante capire che tipo di misurazione si sta andando a fare. Se si misura andando a considerare quella che è la base canonica del circuito, allora il risultato sarà o \(\ket{0}\) o \(\ket{1}\). Non è possibile accedere all'informazione di \(\alpha\) o \(\beta\). Queste sono solo delle probabilità per cui il sistema potrà trovarsi in un autospazio o nell'altro. Se misurassimo con un dispositivo che fa riferimento, invece, alla base di Hadamard, allora i possibili risultati sarebbero \(\ket{+}\) e \(\ket{-}\). 

Quindi lo stato di un qubit collassa sul uno dei vettori ortogonali di riferimento, utilizzati per la misura, con probabilità definite dai coefficenti del qubit stesso.
Dopo la misurazione, il qubit collassato, non porterà ulteriore informazione rispetto al passato, quindi altre misurazioni, dopo la prima, daranno sempre lo stesso risultato. Il Qubit, quindi, "collasserà in un bit".

La misura avviene lungo uno specifico operatore. Gli operatori sono le matrici di Pauli. Sì dice cos', poichè gli autovettori associati alle matrici di pauli, costituiranno la base in cui il qubit potrà collassare.

Questo perchè ad ogni grandezza fisica osservabile, corrisponde un operatore autoaggiunto il cui spettro costituisce l'insieme dei possibili valori che potranno essere ottenuti dalla misura di tale grandezza.

Questo per dire che la superposizione non solo dipende dalla base di interpretazione, ma anche dall'operatore su di cui si andrà ad effettuare la misurazione.

Alcuni esempi di tali stati:
\[
\ket{\varphi} = 1 \ket{0} + 0 \ket{1}
\]
\[
H = \frac{1}{\sqrt{2}}
\begin{bmatrix}
1 & 1 \\
1 & -1
\end{bmatrix}
\]
\[
H \ket{\varphi} = \frac{1}{\sqrt{2}}
\begin{bmatrix}
1 & 1 \\
1 & -1    
\end{bmatrix}
\begin{bmatrix}
1 \\
0
\end{bmatrix}
= \frac{1}{\sqrt{2}}
\begin{bmatrix}
1 \\
1
\end{bmatrix}
= \frac{1}{\sqrt{2}} (\ket{0} + \ket{1})
\]

Se noi misurassimo tale qubit sull'operatore \(Z\), quindi sulla base canonica, allora avremmo che esso potrà collassare con una probabilità del \(50 \%\), sia in \(\ket{0}\) sia in \(\ket{1}\). Mentre se lo misurassimo lungo l'operatore \(X\), quindi con la base di Hadamard, allora,  si otterrebbe direttamente lo stato \(\ket{+}\).

Quindi la superposizione è dipendente dalla base, poichè, con la misurazione in alcune basi, si otterrebbe un risultato deterministico, mentre in altre si otterrebbe un risultato probabilistico (tale affermazione è alla base del funzionamento del sistema QKD).

Quindi, generalizzando anche ai casi multi-qubit. Un sistema multi-qubit, quando misurato, può cadere in una delle sue \(2^n\) possibili vettori della base. Pertanto si formalizza quello che è il terzo postulato della meccanica quantistica:

Data una base ortonormale dello spazio di Hilbert corrispondente al nostro sistema fisico quantistico. Ogni misurazione di base avrà \(d\) possibili uscite, corrispondenti ai vettori della base orto-normale. Se un sistema è nello stato \(\ket{\psi}\) e andiamo a misurarlo lungo la base ortonormale \(\{\ket{\psi_n}\}^d_{n=1}\), la probabilità che si otterrà come uscita il vettore \(n\) è data dalla legge di Born:
\[
p(n|\ket{\psi}) = |\braket{\psi_n | \psi}|^2
\]

Quindi la probabilità non è altro che il modulo quadro della proiezione del Qubit lungo il vettore della base n-esimo. 

Dalla precedente regola, quindi, è possibile andare a definire la somma di tutti i possibili risultati della misurazione e si otterrà la seguente formula:
\[
\sum_{d}^{n=1}p(n | \ket{\psi}) = \sum_{d}^{n=1}|\braket{\psi_n | \psi}|^2 = \braket{\psi | \psi} = 1
\]

Questo dimostra anche che se ci fosse una fase globale, questa non modificherebbe in alcun modo le probabilità associate ai singoli vettori della base. Quindi anche se due stati diversi, hanno tutte le probaiblità uguali, allora essi saranno due stati che differiscono della sola fase globale.

\subsection{L'esperimento di Dirac}
L'unica particella visibile all'umo è il \textbf{fotone} che possiamo osservare in maniera diretta.

Un fotone viene valutato in base alla sua polarizzazione, tale polarizzazione viene valutata con un piano 2-d qualunque, posizionato perpendicolarmente al movimento del fotone. La scelta di cosa andare a vlutare è arbitraria, cosideriamo le polarizzazioni verticali \(\uparrow\) e polarizzazioni orizzontali \(\rightarrow\), abbiamo scelto la nostra base di valutazione.

\begin{warn}
La base di valutazione è quella che è stata scelta da noi, ma ciò non toglie che si potrebbe valutare tale polarizzazione anche utilizzando altre tipologie di basi ortonormali
\end{warn}

Per effettuare l'esperimento di Dirac, andando a scegliere un pool di filtri A (filtro per polarizzazioni verticali \(\uparrow\)), B (filtro per polarizzazioni \(\nearrow\)) ed un filtro C (filtro per polarizzazioni \(\rightarrow\)).
Prima di iniziare ad aggiungere i filtri, si va a posizionare un generatore di fotoni (es. uno schermo o un laser).
L'esperimento si divide in tre principali momenti:
\begin{enumerate}
    \item Posizionamento del filtro A, che filtrerà la parte di polarizzazione \(\uparrow\) tra la sorgente di fotoni e l'osservatore. Questo farà notare un affievolimento della luce emessa
    
    \item Posizionamento del filtro di tipo C, che filtrerà la restante parte del fotone \(\rightarrow\), questo non 

\end{enumerate}